% TOP_PROBLEM: {"id": 34, "title": "The Volume Conjecture", "category_id": 7, "category": {"id": 7, "name": "topology", "display_name": "Topology", "description": "Properties preserved under continuous deformations.", "slug": "topology", "order_index": 7, "created_at": "2026-07-31T15:26:25.671Z"}, "status": "open"}
% FIRST_POSTED: 2026-09-25
% ENTRY_KIND: statement_only
% !TeX program = lualatex
% Definition and sources only; this file is not an LLM solution attempt.
\documentclass[11pt]{article}
\usepackage[margin=25mm]{geometry}
\usepackage{fontspec}
\setmainfont{Latin Modern Roman}
\usepackage{amsmath,amssymb,mathtools}
\usepackage{unicode-math}
\setmathfont{Latin Modern Math}
\usepackage{xurl,hyperref}
\hypersetup{colorlinks=true,urlcolor=blue}
\setcounter{secnumdepth}{0}
\setlength{\parindent}{0pt}
\setlength{\parskip}{5pt}
\setlength{\emergencystretch}{3em}
\begin{document}
\section*{\texorpdfstring{The Volume Conjecture}{The Volume Conjecture}}\label{34}
\subsection{Definitions and mathematical statement}
For a knot $K\subset S^3$, let $J_{K,N}(q)$ be the $N$-colored Jones polynomial normalized so that the unknot has value one. A hyperbolic knot has a complete finite-volume hyperbolic metric on its complement. The conjecture is
\[
 \lim_{N\to\infty}\frac{2\pi}{N}\log\left|J_{K,N}(e^{2\pi i/N})\right|
 =\operatorname{Vol}(S^3\setminus K).
\]
Both the normalization and the root-of-unity convention matter. The statement includes existence of the limit for every hyperbolic knot; a numerical fit or a result for one family is not the universal conclusion.
\par\smallskip\textit{Formulation and concept references:} \href{https://ems.press/content/serial-article-files/51448?nt=1}{[S1]}.

\subsection{Short English statement}
Does the exponential growth of colored Jones values at roots of unity recover the hyperbolic volume of every hyperbolic knot complement?
\subsection{Sources}
\begin{itemize}
\item[S1]C. Aistleitner and B. Borda, Journal of the European Mathematical Society 27 (2025), 5043-5092.
\url{https://ems.press/content/serial-article-files/51448?nt=1}.
\textit{Formulation source.}
\end{itemize}
\end{document}
